\answer{ex:14:1}{(a)~De $0.17$~s et de $1.5$~s. (b)~À une distance de moins de $11$~cm.}
\answer{ex:14:2}{Un toucher sans douleur ne passe pour aucun $\mu$ tant que $g\le\beta w_{q\mu}/w_{z\mu}=5$; un toucher $\mu=1$ passe pour $g>6$: une forte sensibilisation rend douloureux un contact ordinaire.}
\answer{ex:14:3}{(a)~$g^*=(1-k_sC)/(1-k_sA)$, stable pour $k_sA<1$. (b)~Sans toucher, pour $\nu\ge2$; avec toucher, dès $\nu\ge1.7$.}
\answer{ex:14:4}{$V(t)=-Pe^{-(T-t)/\tau_\gamma}$. (a)~$-Pe^{-T/\tau_\gamma}\approx-0.37P$. (b)~$+Pe^{-(T-t_e)/\tau_\gamma}\approx+0.61P$; elle croît avec $t_e$ jusqu'à $+P$ et enseigne à éviter la douleur.}
\answer{ex:14:5}{Pour $k_s=4$, pas de verrouillage (un second état stable existe pour $k_s\ge4.58$); $T_{\min}\approx4.35$, $1.40$, $0.65$, $0.32\,\tau_s$ pour $k_s=4.6$, $5$, $6$, $8$.}
\answer{ex:14:6}{$w_{q\nu}=4/9\approx0.444$, $q_0=2/9\approx0.222$, $w_{q\mu}=49/45\approx1.089$; un toucher sans douleur est sans danger jusqu'à $g=49/9\approx5.4$.}
\answer{ex:14:7}{(a)~Seuil $\nu_c(g)=\theta/g$ si $\nu_c\ge q_0/w_{q\nu}$, sinon $(\theta+\beta q_0)/(g+\beta w_{q\nu})$; $g^*\approx2.45$, $1.0$, $0.25$. (b)~Question ouverte.}
